\section{赌徒破产中 OST 的条件失效}

\subsection{背景介绍}

赌徒破产问题是概率论中最古老的随机过程模型之一，其历史可追溯至 Pascal 和 Fermat 关于赌博问题的通信\cite{feller1968introduction}。从现代鞅论的角度看，它是研究 OST 的天然试验场：过程是简单的对称随机游走（最经典的鞅），停时是首达吸收边界，且可以通过选择单边或双边边界来切换 OST 是否成立。

考虑赌徒甲初始资金为 $a$ 元，赌徒乙初始资金为 $b$ 元。每局甲以概率 $p$ 赢 1 元，以概率 $q = 1-p$ 输 1 元。令 $S_n$ 为甲在前 $n$ 局后的净赢额（$S_0 = 0$），即 $S_n = \sum_{i=1}^{n} \xi_i$，其中 $\xi_i = \pm 1$。当 $p = 1/2$ 时，$\{S_n\}$ 是一个鞅；当 $p \neq 1/2$ 时，$\{S_n\}$ 不是鞅，但可以构造 Wald 鞅 $W_n = (q/p)^{S_n}$ 或补偿鞅 $S_n - (2p-1)n$。

本文第四章聚焦于 $p = 1/2$（即公平赌博）的情况。此时 OST 成立与否完全取决于停时的选择。

\subsection{理论分析}

\subsubsection{情形 A：双边障碍——OST 完美成立}

定义停时 $\tau_{\text{two}} = \inf\{n: S_n = -a\ \text{或}\ S_n = +b\}$（甲输光或乙输光）。由于 $\tau_{\text{two}}$ 是双边首出时，$S_{n \wedge \tau_{\text{two}}}$ 始终被约束在 $[-a, b]$ 之内，故 $|S_{n \wedge \tau_{\text{two}}}| \leqslant \max(a, b)$。OST 的条件 2（被停过程有界）满足。此外，随机游走在有限区间内的首出时具有有限的期望值，$\mathbb{E}[\tau_{\text{two}}] < \infty$。因此 OST 无条件成立，$\mathbb{E}[S_{\tau_{\text{two}}}] = \mathbb{E}[S_0] = 0$，即
\begin{equation}\label{eq:gambler_ost}
  -a \cdot \mathbb{P}(S_{\tau_{\text{two}}} = -a) + b \cdot \mathbb{P}(S_{\tau_{\text{two}}} = +b) = 0.
\end{equation}
解之得甲的输光概率和乙的输光概率分别为
\begin{equation}\label{eq:gambler_prob}
  P_a = \frac{b}{a+b}, \quad P_b = \frac{a}{a+b}.
\end{equation}
这就是经典的赌徒破产公式：甲的破产概率等于乙的初始资金占总资金的比例。

为了求平均吸收时间，再考虑过程 $\{S_n^2-n\}$。记随机游走增量为 $\eta_n=S_n-S_{n-1}$，则 $\eta_n=\pm1$ 且
\[
\mathbb{E}[\eta_n\mid\mathcal{F}_{n-1}] = 0,\qquad \eta_n^2=1.
\]
于是对任意 $n\geqslant1$，
\begin{align}
  \mathbb{E}[S_n^2-n\mid\mathcal{F}_{n-1}]
  &= \mathbb{E}\big[(S_{n-1}+\eta_n)^2-n \mid \mathcal{F}_{n-1}\big] \notag\\
  &= S_{n-1}^2 + 2S_{n-1}\mathbb{E}[\eta_n\mid\mathcal{F}_{n-1}]
     + \mathbb{E}[\eta_n^2\mid\mathcal{F}_{n-1}] - n \notag\\
  &= S_{n-1}^2 - (n-1).
\end{align}
因此 $\{S_n^2-n\}$ 的确是鞅。对有界停时 $\tau_{\text{two}}\wedge m$ 应用 OST，有
\[
\mathbb{E}[S_{\tau_{\text{two}}\wedge m}^2-(\tau_{\text{two}}\wedge m)] = 0.
\]
当 $m\to\infty$ 时，由 $|S_{\tau_{\text{two}}\wedge m}|\leqslant \max(a,b)$ 及控制收敛定理有 $\mathbb{E}[S_{\tau_{\text{two}}\wedge m}^2] \to \mathbb{E}[S_{\tau_{\text{two}}}^2]$；而 $\tau_{\text{two}}\wedge m \uparrow \tau_{\text{two}}$，由单调收敛定理有 $\mathbb{E}[\tau_{\text{two}}\wedge m] \uparrow \mathbb{E}[\tau_{\text{two}}]$。令 $m\to\infty$ 即得
\[
\mathbb{E}[S_{\tau_{\text{two}}}^2] = \mathbb{E}[\tau_{\text{two}}].
\]
再结合 $S_{\tau_{\text{two}}}\in\{-a,b\}$ 与式 \eqref{eq:gambler_prob}，
\begin{equation}\label{eq:gambler_etau}
  \mathbb{E}[\tau_{\text{two}}]
  = a^2 P_a + b^2(1-P_a)
  = a^2 \frac{b}{a+b} + b^2 \frac{a}{a+b}
  = ab.
\end{equation}
例如，甲有 20 元、乙有 10 元时，平均需要 $20 \times 10 = 200$ 局才有一方输光。

\subsubsection{情形 B：单边障碍——OST 的条件失效}

现在定义停时 $\tau_{\text{one}} = \inf\{n: S_n = +b\}$。为使单边首达时具有物理意义，此处允许甲无限负债——无论过程在负半轴游走到多深的位置，博弈均不终止；甲仅在净收益首次触及 $+b$ 时主动停手。一维随机游走的常返性保证 $\tau_{\text{one}} < \infty$ 几乎必然（即游走几乎必然在某有限时刻抵达 $+b$），然而：

\begin{enumerate}[noitemsep]
  \item \textbf{停时无界}：对任意 $K$，$\mathbb{P}(\tau_{\text{one}} > K) > 0$。游走可以在碰到 $+b$ 之前在负半轴逗留任意长的时间。
  \item \textbf{过程在停时前无（下）界}：在 $\tau_{\text{one}}$ 之前，$S_n$ 可以取任意大的负值——游走可以先去到 $-1, -10, -100$，再折返碰到 $+b$。
  \item \textbf{停时期望无限}：$\mathbb{E}[\tau_{\text{one}}] = \infty$，其原因在于停时分布的尾部行为（见下文）。
\end{enumerate}
条件 1、2、3 全部不成立。相应地，OST 失效：
\begin{equation}
  \mathbb{E}[S_{\tau_{\text{one}}}] = b \neq 0 = \mathbb{E}[S_0].
\end{equation}

\textbf{失效的深层原因：一致可积性的缺失。} OST 失效的数学本质在于被停过程 $\{S_{n \wedge \tau_{\text{one}}}\}$ 不是一致可积的\cite{chow1997probability}。回顾一族随机变量 $\{X_\alpha\}$ 一致可积的定义：$\lim_{M \to \infty} \sup_{\alpha} \mathbb{E}[|X_\alpha| \cdot \mathbf{1}_{\{|X_\alpha| > M\}}] = 0$。对于 $\{S_{n \wedge \tau_{\text{one}}}\}$，虽然 $\tau_{\text{one}} < \infty$ 几乎必然，但那些在 $\tau_{\text{one}} > n$ 条件下的路径上 $S_n$ 可能取到非常大的负值，使得对任意大的 $M$，始终存在某个 $n$ 使尾部期望不可忽略。一致可积性的缺失直接导致极限与期望不能交换：
\begin{equation}
  \lim_{n \to \infty} \mathbb{E}[S_{n \wedge \tau_{\text{one}}}] = 0 \quad \neq \quad
  \mathbb{E}\Big[\lim_{n \to \infty} S_{n \wedge \tau_{\text{one}}}\Big] = \mathbb{E}[S_{\tau_{\text{one}}}] = b.
\end{equation}

\textbf{停时分布的尾部行为。} 双边停时和单边停时的尾部行为有本质区别。对于双边障碍 $[-a, b]$，由生灭链的谱理论知 $\mathbb{P}(\tau_{\text{two}} > n)$ 以指数速率衰减，故 $\mathbb{E}[\tau_{\text{two}}] < \infty$。对于单边障碍 $\tau_{\text{one}} = \inf\{n: S_n = +b\}$，由随机游走理论可得（当 $b = 1$ 时）\cite{feller1968introduction}
\begin{equation}\label{eq:tail_heavy}
  \mathbb{P}(\tau_{\text{one}} > n) \sim \frac{\sqrt{2}}{\sqrt{\pi n}}, \quad n \to \infty.
\end{equation}
这是一个以 $n^{-1/2}$ 速率衰减的重尾分布，其积分发散，因此 $\mathbb{E}[\tau_{\text{one}}] = \infty$。重尾性意味着绝大多数路径会在相对较短的时间内碰到 $+b$，但极少数顽固路径（游走先大幅下跌再缓慢回升）贡献了无穷的期望。

\textbf{截断停时与不可交换极限。} 为展示 OST 失效的精确机制，引入截断停时 $\tau_N = \min(\tau_{\text{one}}, N)$。对任意固定的 $N$，$\tau_N$ 有界（$\leqslant N$），OST 的条件 1 满足，因此
\begin{equation}\label{eq:trunc_ost}
  \mathbb{E}[S_{\tau_N}] = \mathbb{E}[S_0] = 0, \quad \forall N < \infty.
\end{equation}
另一方面，由于 $\tau_{\text{one}} < \infty$ 几乎必然成立，随机变量序列逐点收敛：$S_{\tau_N} \xrightarrow[N\to\infty]{a.s.} S_{\tau_{\text{one}}}=b$。于是出现了 OST 反例最核心的现象：
\begin{equation}\label{eq:no_ui_limit}
  \lim_{N\to\infty}\mathbb{E}[S_{\tau_N}] = 0,
  \qquad
  \mathbb{E}\!\left[\lim_{N\to\infty}S_{\tau_N}\right] = b.
\end{equation}
二者不相等，说明 $\{S_{\tau_N}\}_{N\geqslant 1}$ 不是一致可积的。直观地看，在大多数已经命中 $+b$ 的路径上，$S_{\tau_N}=b$；但少数尚未命中的路径可能位于很深的负半轴，其负值贡献恰好抵消前者，使每个固定 $N$ 的期望仍为 0。随着 $N$ 增大，这些负值路径越来越稀少但也越来越极端。

\subsection{数值实验}

设置 $a = 20$、$b = 10$（甲有 20 元初始资金，乙有 10 元），即双边障碍为 $[-20, 10]$。在此非对称条件下，理论预测甲的破产概率为 $b/(a+b) = 1/3$，期望吸收时间为 $ab = 200$ 局。对于尾部实验，用 3000 条路径（最大步数 5000）估计双边和单边停时的生存函数 $\mathbb{P}(\tau > t)$。各实验的具体参数见表~\ref{tab:gambler_params}。

\begin{table}[H]
  \centering
  \caption{赌徒破产模型数值实验条件}
  \label{tab:gambler_params}
  \begin{tabular}{llll}
    \toprule
    参数 & 图 4.1 & 图 4.2 & 图 4.3 \\
    \midrule
    初始资金 $(a,b)$ & $(20,10)$ & $(20,10)$ & $(20,10)$ \\
    胜率 $p$ & 0.5 & 0.5 & 0.5 \\
    双边障碍 & $[-20, 10]$ & — & $[-20, 10]$ \\
    单边障碍 & — & $\tau_{\text{one}} = \inf\{n:S_n=+10\}$ & $\tau_{\text{one}} = \inf\{n:S_n=+10\}$ \\
    路径数 & 12（展示） & 16（展示） & 3000 \\
    最大步数 & 10000 & 10000 & 5000 \\
    实验目标 & 双边样本轨道 & 单边样本轨道与截断 & 停时尾部比较 \\
    \bottomrule
  \end{tabular}
\end{table}

在双边障碍 $[-20, 10]$ 下，对称随机游走在两个吸收壁之间来回游荡（图~\ref{fig:ch04_paths}）。虽然每条路径的最终停时值为 $+10$ 或 $-20$，但按吸收概率加权平均后恰为零——蓝色路径被 $+10$ 吸收、红色路径被 $-20$ 吸收的比例符合理论预测的 $2/3$ 与 $1/3$。这一样本层面的表现与 OST 在双边情形下的严格成立完全一致。

\begin{figure}[H]
  \centering
  \includegraphics[width=0.9\textwidth]{ch04_paths}
  \caption{双边障碍下的样本轨道（$a=20, b=10$）。蓝色路径最终被上界 $+10$ 吸收；红色路径最终被下界 $-20$ 吸收。圆点标记吸收时刻。}
  \label{fig:ch04_paths}
\end{figure}

单边障碍则将系统推入了 OST 的失效区域。图~\ref{fig:ch04_unilateral} 以上下两个子图对照了路径的两种典型命运：绝大多数在数百步内即触及 $+10$（上图），与 $\tau_{\text{one}} < \infty$ 几乎必然的理论结论一致；但仍有少数路径在 10000 步的观察窗口内始终未能达到目标（下图）——它们先向负方向大幅游走，随后始终未能回升至 $+10$。这些极端路径尽管稀少，却正是 OST 失效的微观来源：它们在负半轴累积了任意大的负值，使得被停过程 $\{S_{n \wedge \tau_{\text{one}}}\}$ 丧失了一致可积性，式~\eqref{eq:no_ui_limit} 中极限与期望的交换因而缺乏数学依据。

\begin{figure}[H]
  \centering
  \includegraphics[width=0.9\textwidth]{ch04_unilateral}
  \caption{单边障碍下的样本轨道（$b=10$）。上子图：在最大步数内成功达到 $+10$ 的路径。下子图：在最大步数内仍未达到 $+10$ 的截断路径——这些极端路径先向负方向大幅游走，严重拖慢了回升过程，是 OST 失效的微观来源。}
  \label{fig:ch04_unilateral}
\end{figure}

停时分布的尾部行为从统计层面解释了这一机制的必然性。通过 3000 条独立路径估计的经验生存函数（图~\ref{fig:ch04_tail}）显示，双边停时的尾部呈指数级快速衰减（轻尾），对应 $\mathbb{E}[\tau_{\text{two}}] < \infty$；而单边停时在对数-对数坐标下呈现出与 $C t^{-1/2}$ 平行的缓慢衰减（重尾），其积分发散，故 $\mathbb{E}[\tau_{\text{one}}] = \infty$。重尾性的存在意味着图~\ref{fig:ch04_unilateral} 下子图中那些顽固路径并非偶然的模拟噪声，而是单边随机游走停时分布的固有数学特征——只要路径数足够多、观察窗口足够长，就必然会出现。

\begin{figure}[H]
  \centering
  \includegraphics[width=0.9\textwidth]{ch04_tail}
  \caption{停时生存函数 $\mathbb{P}(\tau > t)$ 的双对数图。3000 条独立路径的经验生存函数：双边停时为轻尾，$\mathbb{E}[\tau] < \infty$；单边停时呈 $C t^{-1/2}$ 型重尾，$\mathbb{E}[\tau] = \infty$。重尾性是 OST 条件 3 失效的微观原因。}
  \label{fig:ch04_tail}
\end{figure}
